\section{D10 Heat-Kernel Calibration Supplement}
\label{app:particle-d10-heat-kernel}

This appendix carries the D10 compact-group heat-kernel items needed by the particle surface.
The theorem-bearing leaf for the same-overlap thermal/Casimir law sits on the
screen-microphysics paper surface at fixed cutoff. What is recorded here is the
compact-group / Peter--Weyl lift and beta-shift bookkeeping used by the particle-side
calibration branch. This claim tier does not derive the physical one-loop beta coefficients
from edge sectors alone.

\subsection{\texorpdfstring{\(\mathcal R_U\)}{R\_U} Unified-Coupling Certificate}

The \(\mathcal R_U\) certificate is the D10 gate needed by the hierarchy discussion in the
synthesis paper. It certifies the unified diffusion coupling used in the electroweak
transmutation factor
\[
\frac{v}{E_\star}
=
P^{-1/2}\exp\!\left[-\frac{2\pi}{4\alpha_U(P)}\right].
\]
It does not certify the full clock branch
\[
\mathcal R_\gamma
=
\mathcal R_U+\mathcal R_\alpha+\mathcal R_e^{\mathrm{abs}}
+\mathcal R_{\mathrm{QCD/nuc}}^{133\mathrm{Cs}}
+\mathcal R_{\mathrm{atom}}^{133\mathrm{Cs}}.
\]
The full \(\mathcal R_\gamma\) stack requires source-only electromagnetic endpoint,
charged-lepton absolute-scale, cesium nuclear, and atomic spectral enclosures. The numerical package is
a numerical witness for \(\mathcal R_U\), with formal promotion gated by outward-rounded interval
arithmetic.

\paragraph{Frozen source packet.}
The public-pixel branch uses
\[
P_C
=
\decnum{1.6309682094039593248792798477826489413359828516279250606661507533907793398933432}.
\]
The source-audit branch uses the certified coordinate
\[
P_{\mathrm{fwd}}
=
1.630972095858897\ldots.
\]
It is interval-certified unique on its interval and, by the domain-global certificate, the only fixed point of its
readout map on the declared physical domain \(\alpha^{-1}\in[100,200]\)
(the interval contraction certificate of 2026-07-14,
\texttt{p\_global\_uniqueness\_certificate\_2026-07-14.json};
closure row CL-6, closed).
The finite heat-kernel cutoffs are
\[
N_2=128,\qquad N_3=64.
\]
The one-loop running packet is
\[
(b_1,b_2,b_3)=\left(\frac{33}{5},1,-3\right),
\qquad
N_c=3,
\qquad
\beta_{\mathrm{EW}}=N_c+1=4.
\]
This packet is frozen in the certificate. It is not silently derived in this appendix from the
recovered OPH core.

\begin{definition}[No-hidden-free-variable rule for \(\mathcal R_U\)]
The following objects must be fixed before comparison with electroweak or gravity data:
\[
(b_1,b_2,b_3),\quad
N_2,\quad
N_3,\quad
M_U(P),\quad
E_{\mathrm{cell}}(P),\quad
\beta_{\mathrm{EW}},\quad
\hbox{Casimir normalization},\quad
\hbox{matching scheme}.
\]
Changing any of them after comparison with \(M_W\), \(M_Z\), \(\alpha_s(M_Z)\), low-energy
electroweak couplings, or \(G_{\mathrm{exp}}\) demotes the row to calibration.
\end{definition}

\paragraph{Dimensionless source map.}
Set \(E_\star=1\) for the dimensionless calculation. For a trial coupling \(a\), define
\[
M_U(P)=e^{-2\pi}P^{1/6},
\qquad
E_{\mathrm{cell}}(P)=P^{-1/2},
\]
\[
v(P,a)=E_{\mathrm{cell}}(P)\exp\!\left[-\frac{2\pi}{4a}\right].
\]
For \(i=1,2,3\),
\[
\alpha_i^{-1}(\mu;P,a)
=
a^{-1}
+
\frac{b_i}{2\pi}\log\!\left(\frac{M_U(P)}{\mu}\right),
\qquad
\alpha_Y(\mu;P,a)=\frac35\alpha_1(\mu;P,a).
\]
The source \(Z\)-scale is the positive solution of
\[
\mu
=
\frac{v(P,a)}{2}
\sqrt{4\pi\alpha_2(\mu;P,a)+4\pi\alpha_Y(\mu;P,a)}.
\]
For \(\mathrm{SU}(2)\), with \(j=n/2\) and \(0\le n\le N_2\),
\[
d_j=2j+1,\qquad C_2(j)=j(j+1).
\]
For \(\mathrm{SU}(3)\), with \(0\le p,q\le N_3\),
\[
d_{p,q}=\frac{(p+1)(q+1)(p+q+2)}{2},
\qquad
C_2(p,q)=\frac{p^2+q^2+pq+3p+3q}{3}.
\]
The group sums are
\[
Z_G(t)=\sum_R d_R e^{-tC_2(R)},
\qquad
\bar\ell_G(t)=
\frac{1}{Z_G(t)}
\sum_R d_R e^{-tC_2(R)}\log d_R.
\]
Define
\[
t_2(P,a)=4\pi^2\alpha_2(\mu_Z(P,a);P,a),
\qquad
t_3(P,a)=4\pi^2\alpha_3(\mu_Z(P,a);P,a),
\]
and
\[
\Phi_U(P,a)
=
\bar\ell_{\mathrm{SU}(2)}(t_2(P,a))
+
\bar\ell_{\mathrm{SU}(3)}(t_3(P,a))
-
\frac{P}{4}.
\]
The certificate value is the zero \(\Phi_U(P,\alpha_U(P))=0\).

\paragraph{Numerical witness.}
For \(P=P_C\), the checker emits
\[
\begin{aligned}
\alpha_U(P_C)&=0.041124336195630495,\\
\alpha_U(P_C)^{-1}&=24.316501918546496.
\end{aligned}
\]
The corresponding dimensionless hierarchy data are
\[
\begin{aligned}
\frac{M_U(P_C)}{E_\star}&=0.0020260720012100037,\\
\frac{v(P_C)}{E_\star}&=2.0199803239725553\times10^{-17},\\
\frac{\mu_Z(P_C)}{E_\star}&=7.502403238986022\times10^{-18}.
\end{aligned}
\]
The emitted couplings at \(\mu_Z\) are
\[
\begin{aligned}
\alpha_1&=0.016885708038988166,\\
\alpha_Y&=0.010131424823392899,\\
\alpha_2&=0.033777889448908055,\\
\alpha_3&=0.11833602747445048.
\end{aligned}
\]
The heat-kernel parameters and entropy readouts are
\[
\begin{aligned}
t_2&=1.3334976254578108,\\
t_3&=4.6717191102770705,
\end{aligned}
\]
\[
\begin{aligned}
\bar\ell_{\mathrm{SU}(2)}&=0.39488144681077636,\\
\bar\ell_{\mathrm{SU}(3)}&=0.012860605540213493.
\end{aligned}
\]
Thus
\[
\begin{aligned}
\bar\ell_{\mathrm{SU}(2)}+\bar\ell_{\mathrm{SU}(3)}
&=0.40774205235098987,\\
P_C/4&=0.4077420523509898,
\end{aligned}
\]
and
\[
\begin{aligned}
\Phi_U(P_C,0.041124336195630495)
&=5.55\times10^{-17}.
\end{aligned}
\]
A centered derivative check gives
\[
\partial_a\Phi_U(P_C,a)\simeq -10.990524683118785.
\]
With \(\delta=10^{-6}\), the bracket signs are
\[
\Phi_U(P_C,\alpha_U-\delta)=1.0990748697592423\times10^{-5},
\]
\[
\Phi_U(P_C,\alpha_U+\delta)=-1.0990300680135956\times10^{-5}.
\]

For the certified source-audit pixel,
\[
\begin{aligned}
P_{\mathrm{fwd}}&=1.630972095858897\ldots,\\
\alpha_U(P_{\mathrm{fwd}})&=0.041124247441816685\ldots,\\
\frac{v(P_{\mathrm{fwd}})}{E_\star}&=2.0198114078576331\times10^{-17}.
\end{aligned}
\]
This branch is the one to cite when avoiding the measured Thomson endpoint in the upstream pixel
solve.

\begin{theorem}[Electroweak unified-coupling numerical witness]\label{thm:particle-ru-witness}
Assume the frozen source map, running packet, heat-kernel cutoff policy, and
no-hidden-free-variable rule above. Let
\[
I_U=[0.041123336195630494,\;0.041125336195630496].
\]
If outward-rounded interval arithmetic verifies
\[
0\in \Phi_U(P_C,I_U),
\qquad
0\notin \partial_a\Phi_U(P_C,I_U),
\]
then there is a unique \(\alpha_U(P_C)\in I_U\) satisfying
\[
\Phi_U(P_C,\alpha_U(P_C))=0.
\]
Moreover, the dependency graph of this zero contains no measured \(G\), no Planck unit formed
using measured \(G\), no measured \(\Lambda\), no measured \(M_Z\), no measured \(M_W\), and no
measured low-energy gauge coupling.
\end{theorem}

\begin{proof}
Continuity follows because the running couplings, self-consistent \(Z\)-scale solution, and
finite heat-kernel sums are continuous on the declared positive interval. The endpoint sign
change gives existence by the intermediate value theorem. The derivative exclusion gives
uniqueness. The dependency statement follows by inspection of the frozen source map: the residual
is built only from \(P_C\), the frozen running packet, the self-consistent \(Z\)-scale equation,
and finite compact-group representation sums.
\end{proof}

\paragraph{Precision policy.}
The exponential map makes downstream gravity displays sensitive to \(\alpha_U\):
\[
\frac{\partial\log G}{\partial \alpha_U}
\simeq
\frac{\pi}{\alpha_U^2}
\approx
1.86\times10^3.
\]
Therefore the number of digits printed for any downstream \(G\) or \(\varepsilon_{\mathrm{Cs}}\)
row may not exceed the interval precision certified by
\[
\mathcal R_U+\mathcal R_\alpha+\mathcal R_e^{\mathrm{abs}}
+\mathcal R_{\mathrm{QCD/nuc}}^{133\mathrm{Cs}}
+\mathcal R_{\mathrm{atom}}^{133\mathrm{Cs}}.
\]
The witness supports an electroweak hierarchy certificate. It does not support a 52-digit
source-only gravity prediction.

\subsection{One-Loop Calibration Frame}

At one loop, if couplings unify at \((M_U,\alpha_U)\),
\[
\alpha_i^{-1}(M_Z)=\alpha_U^{-1}+\frac{b_i}{2\pi}\ln\frac{M_U}{M_Z}.
\]
On the D10 branch these are the branch values \((M_U(P),\alpha_U(P))\) emitted by the forward
D10 solve; they are not inferred by inverse readback from measured low-energy couplings.
Writing \(A_i:=\alpha_i^{-1}(M_Z)\) and \(L:=\ln(M_U/M_Z)\), one finds
\[
L=\frac{2\pi}{b_1-b_2}(A_1-A_2),
\]
and the corresponding consistency relation for the D10 lane is
\[
A_3^{\mathrm{pred}}
=
\frac{b_3-b_2}{b_1-b_2}A_1
+
\frac{b_1-b_3}{b_1-b_2}A_2.
\]
This appendix keeps that algebra on the page as a calibration relation, not as a theorem that OPH
derives a supersymmetric UV spectrum.

\subsection{Benchmark coefficient comparison}

For the standard one-loop benchmark coefficients,
\[
\begin{array}{c|c|c}
\text{Model} & (b_1,b_2,b_3) & \alpha_s(M_Z)^{\mathrm{pred}} \\
\hline
\text{SM} & (41/10,-19/6,-7) & \approx 0.071 \\
\text{MSSM} & (33/5,1,-3) & \approx 0.116 \\
\text{Observed} & & 0.1179\pm 0.0010
\end{array}
\]
At one loop, the printed MSSM-style coefficients reproduce the observed closure far better than
the plain SM coefficients. On this D10 branch this is a benchmark matched by the D10
calibration branch, not a theorem that OPH derives a supersymmetric UV spectrum.

\subsection{Heat-Kernel Sector Law}

Read this subsection as the compact-group merge boundary above the fixed-cutoff microphysics
theorem. MaxEnt plus the bi-invariant compact-group package yield the familiar
\(d_R e^{-tC_2(R)}\) law once the Peter--Weyl lift is supplied; they are not a second,
independent proof of the finite-cutoff same-overlap Casimir branch.

\begin{theorem}[Heat-kernel edge-sector weights]
Under MaxEnt with bi-invariant constraints on a compact Lie group \(G\), the sector probabilities
take heat-kernel form:
\[
p_R(t)\propto d_R\,e^{-tC_2(R)},
\]
where \(d_R=\dim R\) and \(C_2(R)\) is the quadratic Casimir.
\end{theorem}

\begin{lemma}[Bi-invariant operators]
If \(G=\prod_i G_i\) is a compact semisimple Lie group written as a product of compact simple
factors, then any bi-invariant second-order differential operator on \(G\) has the form
\[
D=c_0\mathbf 1-\sum_i c_i\Delta_{G_i},
\]
where \(\Delta_{G_i}\) is the Laplace--Beltrami operator on the factor \(G_i\).
\end{lemma}

\begin{remark}
Bi-invariance is equivalent to \(D\in Z(U(\mathfrak g))\). Linear terms vanish because there are
no invariant vectors in the adjoint representation, and the quadratic part is proportional to the
Killing form on each simple factor, giving one independent coefficient per factor. The Casimir
element then acts as \(-\Delta_G\) in the regular representation.
\end{remark}

\subsection{Peter--Weyl Multiplicity and Beta Shifts}

By Peter--Weyl decomposition, the effective refinement-limit edge representation space is
\[
L^2(G)\cong \bigoplus_R V_R\otimes V_R^*.
\]
Entanglement traces over one factor give multiplicity \(d_R\) in \(p_R\), but loops see both
factors, so the effective multiplicity entering the calibration branch is
\[
N_{\mathrm{eff}}(R)=d_R\,p_R.
\]
The one-loop beta shift from edge sectors is then
\[
\Delta b_a=\sum_R p_R\,d_R\,T_a(R),
\]
where \(T_a(R)\) is the Dynkin index for gauge factor \(a\). This is the exact mathematical point
on the D10 lane where edge-sector heat-kernel weights feed the calibration branch. Matching to
MSSM-like one-loop running is a comparison benchmark, not a theorem-level
MSSM spectrum claim.

The branch boundary for theorem-level unification closure is explicit. One requires a derived
refinement-limit identification of which transportable sectors actually contribute to the running
carrier, a derived statistics/grading rule rather than an imposed fermionic-grading restriction,
controlled threshold and decoupling conventions on that same carrier, a derived representation
content/truncation rule for the sectors retained in the comparison, and a proof that the effective
loop multiplicities entering \(\Delta b_a\) are exactly the ones used in this calibration package
rather than nearby alternatives. The declared runtime surface does not emit the
full RG/matching/threshold/scheme packet: scheme lock, threshold map,
beta provenance, and interval composition are separate certificate requirements.

At the branch unification diffusion parameter \(t_U(P)\approx 1.64\) selected by the forward D10
solve, the benchmark shift is
\[
\Delta b\approx (2.49,4.38,3.97)
\qquad
\text{vs MSSM}
\qquad
(2.50,4.17,4.00).
\]
With the additional D10 calibration assumption of a fermionic-grading restriction to
half-integer \(\mathrm{SU}(2)\) sectors, this sharpens to
\[
\Delta b\approx (2.50,4.17,3.97),
\]
which matches the MSSM-style one-loop benchmark at the percent level under that declared
calibration package. Equivalently, the benchmark ratio
\[
\frac{\Delta b_3}{\Delta b_2}\approx 0.91
\]
sits close to the MSSM comparison value \(0.96\).

\section{H3 Worldline-Stitch Certificate}\label{app:h3-worldline-certificate}

This appendix records the finite evidence-facing certificate for
Theorem~\ref{thm:h3-record-worldline-stitching}. It is separate from the D10 heat-kernel
calibration package above. A passing certificate is a record-continuation statement in a declared
\(H^3\) observer chart, not a particle-species, mass, charge, or scattering-amplitude theorem.

The certificate payload has the following required blocks.
\begin{enumerate}
\item \textbf{Atlas.} It declares the hyperboloid model \(H^3_{R_H}\subset\mathbb R^{1,3}\),
the Lorentz signature \((-+++)\), the curvature radius \(R_H\), chart domains, and chart
transitions \(G_{ji}\in\mathrm{SO}^{+}(1,3)\). Transition residuals must certify Lorentz
inner-product preservation, orientation, and future-sheet preservation.
\item \textbf{Clock map.} Each local record interval is mapped to a common comparison-time line.
The observer-time adjacency margin must dominate the declared clock uncertainty.
\item \textbf{Chart-natural extraction.} Raw local records are connected components of a frozen
detector scalar on the \(H^3\) chart. The detector scalar, thresholds, support collars, and
chart-naturality residuals are included. Implementation record IDs are not inputs.
\item \textbf{Overlap descent.} Records are first descended across genuine chart overlaps. The
payload records same-component join margins, distinct-component separation margins, support
errors, and triple-overlap cocycle residuals.
\item \textbf{Interface crossing.} A candidate stitch must cross a real interface in adjacent
clock intervals. It includes oriented signed-distance margins and a positive lower bound on
normal velocity. Tangential, grazing, and file-boundary-only contacts are rejected.
\item \textbf{Transport.} The candidate edge is evaluated in a common chart after sector and
gauge transport. The atlas holonomy comparison is separate from the gauge-holonomy comparison.
\item \textbf{Assignment.} The matching is one-to-one with appearance/disappearance interval
costs. The proposed winner has an upper cost bound, every competitor has a lower cost bound, and
\[
\Delta_{\mathrm{cert}}
=\min_{M\ne M_*}\underline C(M)-\overline C(M_*)>0 .
\]
Record IDs, global IDs, stitch keys, and shard-local IDs are forbidden in admissibility, costs,
and tie-breaking.
\item \textbf{Refinement.} A coarse/fine pair supplies \(Q_{sr}\), a distortion bound
\(\eta_{sr}\), and a contracted-graph isomorphism check. The coarse gap must satisfy
\(\Delta_r>2\eta_{sr}\).
\item \textbf{Interaction firewall.} A free-propagation slab is required. If an interaction marker
is present, the terminal label is \(\mathrm{H3\_INTERACTION\_REQUIRED}\) and the event is handed
to the interaction solver.
\end{enumerate}

The allowed terminal statuses are
\[
\begin{gathered}
\mathrm{H3\_STITCH\_CERTIFIED},\quad
\mathrm{H3\_STITCH\_AMBIGUOUS},\quad
\mathrm{H3\_STITCH\_REJECTED},\\
\mathrm{H3\_INTERACTION\_REQUIRED},\quad
\mathrm{H3\_ATLAS\_INVALID},\quad
\mathrm{H3\_CERTIFICATE\_INCOMPLETE}.
\end{gathered}
\]
\(\mathrm{H3\_STITCH\_CERTIFIED}\) is emitted only when every block above passes and the robust
assignment gap is positive. Equal-cost matchings, preserved nontrivial permutations of the visible
data, or coarse/fine gaps below margin force \(\mathrm{H3\_STITCH\_AMBIGUOUS}\). Missing fields
force \(\mathrm{H3\_CERTIFICATE\_INCOMPLETE}\). Failed atlas, interface, metric, or transport gates
force rejection or atlas invalidity. In particular, Euclidean distances on
\(\texttt{h3SpatialPoint}\) coordinates are not certificate distances; the certificate distance is
the hyperboloid geodesic \(d_H\) from Section~\ref{sec:h3-record-worldlines}.
